\chapter{Needs across life stages}

\section{Learning objectives}
Relate nutrient requirements to growth, pregnancy, lactation, aging, and training; identify groups at risk of deficiency; and select food-first and supplement strategies appropriate to the life stage.

\section{Clinical principles}
Life-stage recommendations are only useful when affordable, acceptable, and accessible. Assessment should include customary foods, cooking facilities, religious restrictions, disability, food security, and medicines. Growth, pregnancy, lactation, aging, and training change physiology, but they do not justify indiscriminate megadosing.

\section{Pregnancy and breastfeeding}
People who could become pregnant should get 400 mcg of folic acid daily from fortified foods and/or a supplement, starting before conception, to reduce neural-tube-defect risk. Pregnancy increases needs for several nutrients, including folate, iron, iodine, and choline. A prenatal supplement can help meet some needs, but formulations vary and may contain little or no iodine or choline. Review the vitamin A form and dose: excessive preformed vitamin A (retinol or retinyl esters) can harm fetal development; this concern does not apply in the same way to beta-carotene. Supplement choices, including iron and iodine, should be discussed with a qualified clinician. \cite{cdcFolicAcid,odsPregnancy}

Breastfeeding increases some nutrient needs, and maternal diet can affect certain nutrients in breast milk. This does not mean every breastfeeding person needs every supplement; individualize advice, especially for vegan diets or known deficiencies. \cite{cdcMaternalDiet}

\section{Infants and children}
Infants have distinct needs and should not receive adult supplements. In U.S. guidance, breastfed and partially breastfed infants should receive 400 IU of vitamin D daily beginning shortly after birth; a separate supplement is generally not needed when an infant consumes at least 32 ounces of formula per day. Recommendations can differ by country, so follow local pediatric guidance. Iron needs depend on gestational age and feeding: breastfed infants need an iron source by about 6 months, while standard iron-fortified formula generally supplies iron; preterm infants may need additional iron under pediatric supervision. Keep iron, vitamin gummies, and supplements locked away: accidental ingestion can be life-threatening. Growth, appetite, and developmental concerns deserve pediatric assessment rather than empiric megadoses. \cite{cdcInfantVitaminD,cdcInfantIron}

\section{Older adults}
Absorption, appetite, thirst, medication use, and ability to make vitamin D from sunlight can change with age. Stomach changes can also reduce absorption of food-bound vitamin B12; adults over 50 may need to obtain most B12 from fortified foods or supplements, which are more readily absorbed. Attention to protein, vitamin D, calcium, and hydration can also be useful, but routine high-dose supplementation is not automatically beneficial. Fortified foods and appropriately chosen supplements may be easier to tolerate than large tablets. \cite{odsB12,odsD,odsCa}

\section{Vegetarian and vegan patterns}
Well-planned plant-based diets can supply protein, fiber, folate, vitamin C, magnesium, potassium, and many other nutrients. People avoiding animal foods need a reliable vitamin B12 source from fortified foods or a supplement. Iron, zinc, iodine, calcium, and vitamin D sources may need deliberate planning; fortified foods and targeted supplements can help when needed. For omega-3s, plant foods such as flax, chia, and walnuts provide ALA, while algal oil can provide EPA or DHA. \cite{odsB12,odsIron,odsZinc,odsIodine,odsCa,odsD,odsOmega3}

\section{Athletes and physically active people}
Training increases energy and fluid needs, but not necessarily every micronutrient requirement. Sweat sodium losses vary with climate, acclimatization, body size, pace, and individual sweat composition. For prolonged endurance activity, plan fluid replacement around the activity, environment, and individual sweat losses. Overdrinking can cause exercise-associated hyponatremia; electrolyte products do not make excessive fluid intake safe. Avoid rigid one-size-fits-all fluid targets and seek sports-medicine or dietitian guidance for demanding training or competition. \cite{nataFluid}
Sweat losses, energy expenditure, and restrictive eating can raise risk of inadequate intake. A balanced diet and adequate fluids usually matter more than a large micronutrient stack. Recurrent fatigue, injuries, menstrual disruption, or reduced performance should prompt assessment for inadequate energy, iron deficiency, illness, and training load.
